\documentclass[11pt,a4paper]{article}
\usepackage[utf8]{inputenc}
\usepackage[T1]{fontenc}
\usepackage[spanish,es-nodecimaldot]{babel}
\usepackage{amssymb,latexsym,amsmath}
\usepackage[a4paper,top=2.5cm,bottom=2.5cm,left=2.5cm,right=2.5cm,marginparwidth=1.75cm]{geometry}
\usepackage{graphicx}
\usepackage{float}
\usepackage{booktabs}
\usepackage{caption}
\usepackage{subcaption}
\usepackage{enumitem}
\usepackage{listings}
\usepackage{xcolor}
\usepackage{xurl}
\usepackage[colorlinks=true, linkcolor=black, citecolor=black, urlcolor=blue]{hyperref}

% ---------- Configuración de TikZ para diagramas UML nativos ----------
\usepackage{tikz}
\usetikzlibrary{babel,shapes.multipart,arrows.meta,positioning,calc}

\newcommand{\umlm}{$-$\,}
\newcommand{\umlp}{$+$\,}
\newcommand{\estereo}[1]{\guillemotleft\,#1\,\guillemotright}

\tikzset{
  uml/.style={
    rectangle split, rectangle split parts=3,
    draw, fill=white, font=\footnotesize, align=left,
    text width=7.2cm,
    rectangle split part align={center,left,left},
    inner xsep=6pt, inner ysep=4pt
  },
  umltwo/.style={
    uml, rectangle split parts=2,
    rectangle split part align={center,left}
  }
}

% Configuración de colores para código fuente
\definecolor{codegreen}{rgb}{0,0.5,0}
\definecolor{codegray}{rgb}{0.45,0.45,0.45}
\definecolor{codepurple}{rgb}{0.55,0,0.75}
\definecolor{backcolour}{rgb}{0.97,0.97,0.98}

% Configuración de listings para Python y soporte completo UTF-8
\lstset{
    language=Python,
    backgroundcolor=\color{backcolour},
    commentstyle=\color{codegreen}\itshape,
    keywordstyle=\color{blue}\bfseries,
    numberstyle=\tiny\color{codegray},
    stringstyle=\color{codepurple},
    basicstyle=\ttfamily\footnotesize,
    breakatwhitespace=false,
    breaklines=true,
    captionpos=b,
    keepspaces=true,
    numbers=left,
    numbersep=7pt,
    showspaces=false,
    showstringspaces=false,
    showtabs=false,
    tabsize=4,
    frame=single,
    rulecolor=\color{gray!30},
    inputencoding=utf8,
    extendedchars=true,
    literate=
        {á}{{\'a}}1 {é}{{\'e}}1 {í}{{\'i}}1 {ó}{{\'o}}1 {ú}{{\'u}}1
        {Á}{{\'A}}1 {É}{{\'E}}1 {Í}{{\'I}}1 {Ó}{{\'O}}1 {Ú}{{\'U}}1
        {ñ}{{\~n}}1 {Ñ}{{\~N}}1 {¿}{{?`}}1 {¡}{{!`}}1
        {³}{{$^3$}}1 {²}{{$^2$}}1
}

\renewcommand{\contentsname}{Tabla de Contenido}
\renewcommand{\listfigurename}{Índice de Figuras}
\renewcommand{\lstlistingname}{Código}

\begin{document}

% =========================================================================
% PORTADA INSTITUCIONAL UNAL
% =========================================================================
\begin{titlepage}
\center
\begin{minipage}{13.5cm}
\center
\IfFileExists{logo_unal.png}{\includegraphics[width=3.2cm,height=4.1cm]{logo_unal.png}\\[0.4cm]}{\vspace*{1cm}}
\textsc{\Large UNIVERSIDAD NACIONAL DE COLOMBIA}\\[0.4cm]
{\large \textbf{Facultad de Minas}}\\[0.2cm]
{\large \textbf{Sede Medellín}}\\[0.2cm]
{\normalsize Departamento de Ciencias de la Computación y la Decisión}\\[1.3cm]

\rule{\linewidth}{0.5mm}\\[0.4cm]
{\huge \textbf{Programación Orientada a Objetos}}\\[0.3cm]
{\Large \textit{Actividad 2: Clases, Objetos, Métodos y Diagramas UML}}\\[0.2cm]
\rule{\linewidth}{0.5mm}\\[1.3cm]

\begin{tabular}{rl}
    \textbf{Asignatura:} & Programación Orientada a Objetos \\[0.25cm]
    \textbf{Semestre:}   & 2026-2S \\[0.25cm]
    \textbf{Actividad:}  & Actividad 2 -- Clases y Objetos (20\%) \\[0.25cm]
    \textbf{Estudiante:} & Yuliana Stefanía Guzmán Montoya \\[0.25cm]
    \textbf{Cédula:}     & C.C. 1.001.765.162 \\[0.25cm]
    \textbf{Docente:}    & Walter Hugo Arboleda Mazo \\[0.25cm]
    \textbf{Fecha:}      & 8 de octubre de 2026 \\[0.25cm]
\end{tabular}\\[1.5cm]

\textbf{Repositorio GitHub oficial:}\\
\url{https://github.com/yguzmanm-sys/POO-2026-2S}

\vfill
{\large Medellín, Colombia \\ Semestre 2026-2S}
\end{minipage}
\end{titlepage}

\newpage
\tableofcontents
\newpage

% =========================================================================
% INTRODUCCIÓN Y DESCRIPCIÓN GENERAL
% =========================================================================
\section{Descripción General de la Actividad}

El presente documento consolida la solución y el modelado conceptual de los ejercicios correspondientes a la \textbf{Actividad 2} de la asignatura \textit{Programación Orientada a Objetos (2026-2S)}, a partir del libro guía \textit{Ejercicios de programación orientada a objetos con Java y UML} del autor Leonardo Bermón Angarita.

Cada uno de los ejercicios fue implementado en lenguaje \textbf{Python}, adaptando rigurosamente las especificaciones de clases, atributos, visibilidad, constructores, métodos y tipos enumerados. Asimismo, para cada ejercicio se incluye:
\begin{enumerate}[leftmargin=*]
    \item \textbf{Enunciado y Requerimientos:} Descripción formal del problema planteado en el texto guía.
    \item \textbf{URL del Repositorio GitHub:} Enlace directo y verificable al código fuente del ejercicio alojado en GitHub.
    \item \textbf{Código Fuente en Python:} Implementación completa y comentada con tipado estático opcional (\textit{type hints}).
    \item \textbf{Diagrama de Clases UML Nativo en TikZ:} Representación vectorial generada directamente durante la compilación en Overleaf, modelando clases, atributos, métodos y relaciones asociadas sin depender de imágenes externas.
    \item \textbf{Salida de Ejecución:} Resultados obtenidos al invocar el método principal (\texttt{main}), demostrando la validez del funcionamiento.
\end{enumerate}

\newpage

% =========================================================================
% EJERCICIO 2.1
% =========================================================================
\section{Ejercicio 2.1 — Definición de Clases: Clase Persona}
\textbf{Referencia en libro:} Capítulo 2, página 63.

\subsection{Enunciado del problema}
Se requiere un programa que modele el concepto de una persona. Una persona posee nombre, apellido, número de documento de identidad y año de nacimiento. La clase debe tener un constructor que inicialice los valores de sus respectivos atributos.

La clase debe incluir los siguientes métodos:
\begin{itemize}
    \item Definir un método que imprima en pantalla los valores de los atributos del objeto.
    \item En un método \texttt{main} se deben crear dos personas y mostrar los valores de sus atributos en pantalla.
\end{itemize}

\subsection{URL del repositorio en GitHub}
\noindent \textbf{Enlace:} \url{https://github.com/yguzmanm-sys/POO-2026-2S/blob/main/actividad%202/Actividad2_Ejercicio2_1.py}

\subsection{Código fuente en Python}
\begin{lstlisting}[caption={Implementación de la clase Persona (\texttt{Actividad2\_Ejercicio2\_1.py})}]
class Persona:
    def __init__(self, nombre: str, apellido: str, documento: int, nacimiento: int):
        # Constructor que inicializa los atributos
        self.nombre = nombre
        self.apellido = apellido
        self.documento = documento
        self.nacimiento = nacimiento

    def imprimir(self):
        # Metodo que imprime los valores de los atributos
        print(f"Nombre: {self.nombre}")
        print(f"Apellido: {self.apellido}")
        print(f"Documento: {self.documento}")
        print(f"Año de Nacimiento: {self.nacimiento}")

def main():
    # Crear dos personas
    persona1 = Persona("Pedro", "Pérez", 1053121010, 1998)
    persona2 = Persona("Luis", "León", 1053223344, 2001)

    # Imprimir los valores de los atributos de las personas
    print("PERSONA #1")
    persona1.imprimir()
    print()  # Espacio entre las personas
    print("PERSONA #2")
    persona2.imprimir()

# Llamada a la funcion main
if __name__ == "__main__":
    main()
\end{lstlisting}

\subsection{Diagrama de Clases UML}
En el diagrama de clases se define la clase \texttt{Persona} con cuatro atributos (\texttt{nombre}, \texttt{apellido}, \texttt{documento}, \texttt{nacimiento}) y sus métodos (\texttt{\_\_init\_\_} e \texttt{imprimir}). El diagrama es generado de forma vectorial mediante TikZ:

\begin{figure}[H]
    \centering
    \begin{tikzpicture}
        \node[uml] (persona) {\textbf{Persona}
        \nodepart{second}\umlm nombre: str\\
        \umlm apellido: str\\
        \umlm documento: int\\
        \umlm nacimiento: int
        \nodepart{third}\umlp \_\_init\_\_(nombre: str, apellido: str, documento: int, nacimiento: int)\\
        \umlp imprimir(): None};
    \end{tikzpicture}
    \caption{Diagrama de clases UML vectorial para el Ejercicio 2.1 (Clase Persona).}
    \label{fig:uml_ej2_1}
\end{figure}

\subsection{Resultado de ejecución}
\begin{verbatim}
PERSONA #1
Nombre: Pedro
Apellido: Pérez
Documento: 1053121010
Año de Nacimiento: 1998

PERSONA #2
Nombre: Luis
Apellido: León
Documento: 1053223344
Año de Nacimiento: 2001
\end{verbatim}

\newpage

% =========================================================================
% EJERCICIO 2.2
% =========================================================================
\section{Ejercicio 2.2 — Atributos Primitivos y Enumerados: Clase Planeta}
\textbf{Referencia en libro:} Capítulo 2, página 66 (y resolución en páginas 67--72).

\subsection{Enunciado del problema}
Se requiere un programa que modele el concepto de un planeta del sistema solar. Un planeta tiene los siguientes atributos:
\begin{itemize}
    \item Nombre de tipo \texttt{str} con valor inicial \texttt{None}.
    \item Cantidad de satélites de tipo \texttt{int} con valor inicial $0$.
    \item Masa en kilogramos de tipo \texttt{float} con valor inicial $0.0$.
    \item Volumen en kilómetros cúbicos de tipo \texttt{float} con valor inicial $0.0$.
    \item Diámetro en kilómetros de tipo \texttt{int} con valor inicial $0$.
    \item Distancia media al Sol en kilómetros de tipo \texttt{int} con valor inicial $0$.
    \item Tipo de planeta (clasificación por tamaño), mediante un enumerado con los valores: \texttt{GASEOSO}, \texttt{TERRESTRE} y \texttt{ENANO}.
    \item Observable a simple vista, de tipo booleano con valor inicial \texttt{False}.
\end{itemize}

La clase debe incluir los métodos:
\begin{itemize}
    \item Constructor que inicialice los atributos del planeta.
    \item Método para imprimir en pantalla los datos del planeta.
    \item Método para calcular la densidad del planeta como el cociente $\text{masa} / \text{volumen}$.
    \item Método para determinar si es un planeta exterior: un planeta exterior se encuentra más allá de $3.4\text{ UA}$, donde $1\text{ UA} = 149\,597\,870\text{ km}$.
    \item Método \texttt{main} que instancie la Tierra y Júpiter, muestre su información, densidades y si son exteriores.
\end{itemize}

\subsection{URL del repositorio en GitHub}
\noindent \textbf{Enlace:} \url{https://github.com/yguzmanm-sys/POO-2026-2S/blob/main/actividad%202/Actividad2_Ejercicio2_2.py}

\subsection{Código fuente en Python}
\begin{lstlisting}[caption={Implementación de la clase Planeta y enumeración TipoPlaneta (\texttt{Actividad2\_Ejercicio2\_2.py})}]
from enum import Enum

class TipoPlaneta(Enum):
    GASEOSO = 1
    TERRESTRE = 2
    ENANO = 3

class Planeta:
    def __init__(self, nombre: str = None, num_satelites: int = 0, masa: float = 0.0,
                 volumen: float = 0.0, diametro: int = 0, distancia_sol: int = 0,
                 tipo: TipoPlaneta = None, observable: bool = False):
        # Constructor que inicializa los atributos
        self.nombre = nombre
        self.num_satelites = num_satelites
        self.masa = masa
        self.volumen = volumen
        self.diametro = diametro
        self.distancia_sol = distancia_sol
        self.tipo = tipo
        self.observable = observable

    def mostrar_info(self):
        # Metodo que imprime los valores de los atributos
        print(f"Nombre del planeta = {self.nombre}")
        print(f"Cantidad de satélites = {self.num_satelites}")
        print(f"Masa del planeta (kg) = {self.masa}")
        print(f"Volumen del planeta (km³) = {self.volumen}")
        print(f"Diámetro del planeta (km) = {self.diametro}")
        print(f"Distancia media al sol (km) = {self.distancia_sol}")
        print(f"Tipo de planeta = {self.tipo.name if self.tipo else 'No definido'}")
        print(f"Es observable a simple vista = {self.observable}")

    def densidad(self):
        # Calculo de la densidad
        if self.volumen > 0:
            return self.masa / self.volumen
        return 0.0

    def es_exterior(self):
        UA = 149597870
        # Es exterior si esta mas alla de 3.4 UA
        return self.distancia_sol > 3.4 * UA

if __name__ == "__main__":
    # Crear dos planetas
    planeta1 = Planeta("Tierra", 1, 5.9736e24, 1.08321e12, 12742,
                       150000000, TipoPlaneta.TERRESTRE, True)
    planeta2 = Planeta("Júpiter", 79, 1.898e27, 1.4313e15,
                       139820, 778500000, TipoPlaneta.GASEOSO, True)

    # Mostrar info
    planeta1.mostrar_info()
    print(f"Densidad de {planeta1.nombre} = {planeta1.densidad()} kg/km³")
    print(f"¿Es exterior? {planeta1.es_exterior()}\n")

    planeta2.mostrar_info()
    print(f"Densidad de {planeta2.nombre} = {planeta2.densidad()} kg/km³")
    print(f"¿Es exterior? {planeta2.es_exterior()}")
\end{lstlisting}

\subsection{Diagrama de Clases UML}
El diagrama modela la clase \texttt{Planeta} y la clase enumeración \texttt{TipoPlaneta} con el estereotipo \texttt{\estereo{enumeration}}, vinculado mediante una asociación dirigida:

\begin{figure}[H]
    \centering
    \begin{tikzpicture}
        \node[uml, text width=8.5cm] (planeta) {\textbf{Planeta}
        \nodepart{second}\umlm nombre: str = None\\
        \umlm num\_satelites: int = 0\\
        \umlm masa: float = 0.0\\
        \umlm volumen: float = 0.0\\
        \umlm diametro: int = 0\\
        \umlm distancia\_sol: int = 0\\
        \umlm tipo: TipoPlaneta\\
        \umlm observable: bool = False
        \nodepart{third}\umlp \_\_init\_\_(nombre, num\_satelites, masa, volumen, diametro, distancia\_sol, tipo, observable)\\
        \umlp mostrar\_info(): None\\
        \umlp densidad(): float\\
        \umlp es\_exterior(): bool};

        \node[umltwo, text width=4cm, below=1.3cm of planeta] (tipoplaneta) {\estereo{enumeration}\\\textbf{TipoPlaneta}
        \nodepart{second}GASEOSO = 1\\
        TERRESTRE = 2\\
        ENANO = 3};

        \draw[-{Latex[length=2.5mm]}] (planeta.south) -- node[right, font=\scriptsize] {tipo} (tipoplaneta.north);
    \end{tikzpicture}
    \caption{Diagrama de clases UML vectorial para el Ejercicio 2.2 (Clase Planeta y Enum TipoPlaneta).}
    \label{fig:uml_ej2_2}
\end{figure}

\subsection{Resultado de ejecución}
\begin{verbatim}
Nombre del planeta = Tierra
Cantidad de satélites = 1
Masa del planeta (kg) = 5.9736e+24
Volumen del planeta (km³) = 1083210000000.0
Diámetro del planeta (km) = 12742
Distancia media al sol (km) = 150000000
Tipo de planeta = TERRESTRE
Es observable a simple vista = True
Densidad de Tierra = 5514720137369.484 kg/km³
¿Es exterior? False

Nombre del planeta = Júpiter
Cantidad de satélites = 79
Masa del planeta (kg) = 1.898e+27
Volumen del planeta (km³) = 1431300000000000.0
Diámetro del planeta (km) = 139820
Distancia media al sol (km) = 778500000
Tipo de planeta = GASEOSO
Es observable a simple vista = True
Densidad de Júpiter = 1326067211625.7947 kg/km³
¿Es exterior? True
\end{verbatim}

\newpage

% =========================================================================
% EJERCICIO 2.3
% =========================================================================
\section{Ejercicio 2.3 — Estado de un Objeto, Getters y Setters: Clase Automóvil}
\textbf{Referencia en libro:} Capítulo 2, página 73 (y resolución en páginas 74--85).

\subsection{Enunciado del problema}
Se requiere un programa que modele el concepto de un automóvil con los siguientes atributos:
\begin{itemize}
    \item Marca (\texttt{str}), modelo (año de fabricación, \texttt{int}), motor (cilindraje en litros, \texttt{int}).
    \item Tipo de combustible: enumerado (\texttt{GASOLINA}, \texttt{BIOETANOL}, \texttt{DIESEL}, \texttt{BIODIESEL}, \texttt{GAS\_NATURAL}).
    \item Tipo de automóvil: enumerado (\texttt{CARRO\_CIUDAD}, \texttt{SUBCOMPACTO}, \texttt{COMPACTO}, \texttt{FAMILIAR}, \texttt{EJECUTIVO}, \texttt{SUV}).
    \item Número de puertas (\texttt{int}), cantidad de asientos (\texttt{int}), velocidad máxima (\texttt{int}), color (enumerado: \texttt{BLANCO}, \texttt{NEGRO}, \texttt{ROJO}, \texttt{NARANJA}, \texttt{AMARILLO}, \texttt{VERDE}, \texttt{AZUL}, \texttt{VIOLETA}) y velocidad actual (\texttt{int}, inicialmente $0$).
\end{itemize}

La clase debe implementar:
\begin{itemize}
    \item Métodos \textit{getters} y \textit{setters} para todos sus atributos.
    \item Métodos para acelerar, desacelerar y frenar, garantizando que no se supere la velocidad máxima ni se decremente por debajo de $0\text{ km/h}$.
    \item Método para estimar el tiempo de llegada dada una distancia en kilómetros.
    \item Clase de prueba \texttt{Prueba\_Automovil} con método \texttt{main} para simular aceleraciones y frenados controlados.
\end{itemize}

\subsection{URL del repositorio en GitHub}
\noindent \textbf{Enlace:} \url{https://github.com/yguzmanm-sys/POO-2026-2S/blob/main/actividad%202/Actividad2_Ejercicio2_3_Automovil.py}

\subsection{Código fuente en Python}
\begin{lstlisting}[caption={Implementación de la clase Automóvil y pruebas (\texttt{Actividad2\_Ejercicio2\_3\_Automovil.py})}]
from enum import Enum

class TipoCombustible(Enum):
    GASOLINA = 1
    BIOETANOL = 2
    DIESEL = 3
    BIODIESEL = 4
    GAS_NATURAL = 5

class TipoAutomovil(Enum):
    CARRO_CIUDAD = 1
    SUBCOMPACTO = 2
    COMPACTO = 3
    FAMILIAR = 4
    EJECUTIVO = 5
    SUV = 6

class TipoColor(Enum):
    BLANCO = 1
    NEGRO = 2
    ROJO = 3
    NARANJA = 4
    AMARILLO = 5
    VERDE = 6
    AZUL = 7
    VIOLETA = 8

class Automovil:
    def __init__(self, marca: str, modelo: int, motor: int,
                 tipo_combustible: TipoCombustible, tipo_automovil: TipoAutomovil,
                 numero_puertas: int, cantidad_asientos: int,
                 velocidad_maxima: int, color: TipoColor):
        self.marca = marca
        self.modelo = modelo
        self.motor = motor
        self.tipo_combustible = tipo_combustible
        self.tipo_automovil = tipo_automovil
        self.numero_puertas = numero_puertas
        self.cantidad_asientos = cantidad_asientos
        self.velocidad_maxima = velocidad_maxima
        self.color = color
        self.velocidad_actual = 0

    # Getters
    def get_marca(self): return self.marca
    def get_modelo(self): return self.modelo
    def get_motor(self): return self.motor
    def get_tipo_combustible(self): return self.tipo_combustible
    def get_tipo_automovil(self): return self.tipo_automovil
    def get_numero_puertas(self): return self.numero_puertas
    def get_cantidad_asientos(self): return self.cantidad_asientos
    def get_velocidad_maxima(self): return self.velocidad_maxima
    def get_color(self): return self.color
    def get_velocidad_actual(self): return self.velocidad_actual

    # Setters
    def set_marca(self, marca): self.marca = marca
    def set_modelo(self, modelo): self.modelo = modelo
    def set_motor(self, motor): self.motor = motor
    def set_tipo_combustible(self, tipo_combustible): self.tipo_combustible = tipo_combustible
    def set_tipo_automovil(self, tipo_automovil): self.tipo_automovil = tipo_automovil
    def set_numero_puertas(self, numero_puertas): self.numero_puertas = numero_puertas
    def set_cantidad_asientos(self, cantidad_asientos): self.cantidad_asientos = cantidad_asientos
    def set_velocidad_maxima(self, velocidad_maxima): self.velocidad_maxima = velocidad_maxima
    def set_color(self, color): self.color = color
    def set_velocidad_actual(self, velocidad_actual): self.velocidad_actual = velocidad_actual

    def acelerar(self, aceleracion):
        if self.velocidad_actual + aceleracion > self.velocidad_maxima:
            self.velocidad_actual = self.velocidad_maxima
            print("No se puede superar la velocidad máxima.")
        else:
            self.velocidad_actual += aceleracion

    def desacelerar(self, desaceleracion):
        if self.velocidad_actual - desaceleracion < 0:
            self.velocidad_actual = 0
            print("No se puede decrementar a una velocidad negativa.")
        else:
            self.velocidad_actual -= desaceleracion

    def frenar(self):
        self.velocidad_actual = 0

    def calcular_tiempo_llegada(self, distancia):
        if self.velocidad_actual > 0:
            return distancia / self.velocidad_actual
        return float('inf')

    def imprimir(self):
        print(f"Marca = {self.marca}")
        print(f"Modelo = {self.modelo}")
        print(f"Motor = {self.motor}")
        print(f"Tipo de combustible = {self.tipo_combustible.name}")
        print(f"Tipo de automóvil = {self.tipo_automovil.name}")
        print(f"Número de puertas = {self.numero_puertas}")
        print(f"Cantidad de asientos = {self.cantidad_asientos}")
        print(f"Velocidad máxima = {self.velocidad_maxima}")
        print(f"Color = {self.color.name}")

    def imprimir_velocidad_actual(self):
        print(f"Velocidad actual = {self.velocidad_actual}")

class Prueba_Automovil:
    @staticmethod
    def main():
        auto1 = Automovil("Ford", 2018, 3, TipoCombustible.DIESEL,
                          TipoAutomovil.EJECUTIVO, 5, 6, 250, TipoColor.NEGRO)
        auto1.imprimir()
        auto1.set_velocidad_actual(100)
        auto1.imprimir_velocidad_actual()
        auto1.acelerar(20)
        auto1.imprimir_velocidad_actual()
        auto1.desacelerar(50)
        auto1.imprimir_velocidad_actual()
        auto1.frenar()
        auto1.imprimir_velocidad_actual()
        auto1.desacelerar(20)

if __name__ == "__main__":
    Prueba_Automovil.main()
\end{lstlisting}

\subsection{Diagrama de Clases UML}
El diagrama modela la clase \texttt{Automovil}, sus tres enumeraciones asociadas (\texttt{TipoCombustible}, \texttt{TipoAutomovil}, \texttt{TipoColor}) y la clase de prueba \texttt{Prueba\_Automovil}:

\begin{figure}[H]
    \centering
    \resizebox{0.95\textwidth}{!}{
    \begin{tikzpicture}
        \node[uml, text width=8.8cm] (auto) {\textbf{Automovil}
        \nodepart{second}\umlm marca: str\\
        \umlm modelo: int\\
        \umlm motor: int\\
        \umlm tipo\_combustible: TipoCombustible\\
        \umlm tipo\_automovil: TipoAutomovil\\
        \umlm numero\_puertas: int\\
        \umlm cantidad\_asientos: int\\
        \umlm velocidad\_maxima: int\\
        \umlm color: TipoColor\\
        \umlm velocidad\_actual: int = 0
        \nodepart{third}\umlp \_\_init\_\_(marca, modelo, motor, tipo\_combustible, tipo\_automovil, numero\_puertas, cantidad\_asientos, velocidad\_maxima, color)\\
        \umlp get\_marca() / set\_marca(marca)\\
        \umlp get\_modelo() / set\_modelo(modelo)\\
        \umlp get\_motor() / set\_motor(motor)\\
        \umlp get\_tipo\_combustible() / set\_tipo\_combustible(tipo\_combustible)\\
        \umlp get\_tipo\_automovil() / set\_tipo\_automovil(tipo\_automovil)\\
        \umlp get\_numero\_puertas() / set\_numero\_puertas(numero\_puertas)\\
        \umlp get\_cantidad\_asientos() / set\_cantidad\_asientos(cantidad\_asientos)\\
        \umlp get\_velocidad\_maxima() / set\_velocidad\_maxima(velocidad\_maxima)\\
        \umlp get\_color() / set\_color(color)\\
        \umlp get\_velocidad\_actual() / set\_velocidad\_actual(velocidad\_actual)\\
        \umlp acelerar(aceleracion: int): None\\
        \umlp desacelerar(desaceleracion: int): None\\
        \umlp frenar(): None\\
        \umlp calcular\_tiempo\_llegada(distancia: float): float\\
        \umlp imprimir(): None\\
        \umlp imprimir\_velocidad\_actual(): None};

        \node[umltwo, text width=4.2cm, anchor=north west] (comb) at ([xshift=1.8cm]auto.north east) {\estereo{enumeration}\\\textbf{TipoCombustible}
        \nodepart{second}GASOLINA = 1\\
        BIOETANOL = 2\\
        DIESEL = 3\\
        BIODIESEL = 4\\
        GAS\_NATURAL = 5};

        \node[umltwo, text width=4.2cm, anchor=north west] (tauto) at ([yshift=-0.7cm]comb.south west) {\estereo{enumeration}\\\textbf{TipoAutomovil}
        \nodepart{second}CARRO\_CIUDAD = 1\\
        SUBCOMPACTO = 2\\
        COMPACTO = 3\\
        FAMILIAR = 4\\
        EJECUTIVO = 5\\
        SUV = 6};

        \node[umltwo, text width=4.2cm, anchor=north west] (tcolor) at ([yshift=-0.7cm]tauto.south west) {\estereo{enumeration}\\\textbf{TipoColor}
        \nodepart{second}BLANCO = 1\\
        NEGRO = 2\\
        ROJO = 3\\
        NARANJA = 4\\
        AMARILLO = 5\\
        VERDE = 6\\
        AZUL = 7\\
        VIOLETA = 8};

        \node[umltwo, text width=6cm, below=1.4cm of auto] (prueba) {\textbf{Prueba\_Automovil}
        \nodepart{second}\umlp main(): None \quad \{static\}};

        \draw[-{Latex[length=2.5mm]}] (auto.east) -- node[above, sloped, font=\scriptsize, fill=white, inner sep=1pt] {tipo\_combustible} (comb.west);
        \draw[-{Latex[length=2.5mm]}] (auto.east) -- node[above, sloped, font=\scriptsize, fill=white, inner sep=1pt] {tipo\_automovil} (tauto.west);
        \draw[-{Latex[length=2.5mm]}] (auto.east) -- node[above, sloped, font=\scriptsize, fill=white, inner sep=1pt] {color} (tcolor.west);
        \draw[dashed, -{Latex[length=2.5mm]}] (prueba.north) -- node[right, font=\scriptsize] {usa} (auto.south);
    \end{tikzpicture}
    }
    \caption{Diagrama de clases UML vectorial para el Ejercicio 2.3 (Clase Automóvil y Enumeraciones).}
    \label{fig:uml_ej2_3}
\end{figure}

\subsection{Resultado de ejecución}
\begin{verbatim}
Marca = Ford
Modelo = 2018
Motor = 3
Tipo de combustible = DIESEL
Tipo de automóvil = EJECUTIVO
Número de puertas = 5
Cantidad de asientos = 6
Velocidad máxima = 250
Color = NEGRO
Velocidad actual = 100
Velocidad actual = 120
Velocidad actual = 70
Velocidad actual = 0
No se puede decrementar a una velocidad negativa.
\end{verbatim}

\newpage

% =========================================================================
% EJERCICIO 2.4
% =========================================================================
\section{Ejercicio 2.4 — Figuras Geométricas: Círculo, Rectángulo, Cuadrado y Triángulo Rectángulo}
\textbf{Referencia en libro:} Capítulo 2, página 86 (y resolución en páginas 87--94).

\subsection{Enunciado del problema}
Se requiere un programa que modele varias figuras geométricas:
\begin{itemize}
    \item \textbf{Círculo:} radio en centímetros.
    \item \textbf{Rectángulo:} base y altura en centímetros.
    \item \textbf{Cuadrado:} longitud del lado en centímetros.
    \item \textbf{Triángulo Rectángulo:} base y altura en centímetros.
\end{itemize}

Para cada figura se deben definir métodos para calcular su \textbf{área} y su \textbf{perímetro}. Adicionalmente, para el triángulo rectángulo se requiere:
\begin{itemize}
    \item Calcular la longitud de la hipotenusa mediante el teorema de Pitágoras ($h = \sqrt{a^2 + b^2}$).
    \item Determinar el tipo de triángulo según sus lados: equilátero (3 lados iguales), isósceles (2 lados iguales) o escaleno (3 lados diferentes).
\end{itemize}
Se debe desarrollar una clase de prueba \texttt{Figuras\_geometricas} con un método \texttt{main} que cree instancias de las cuatro figuras y muestre los cálculos.

\subsection{URL del repositorio en GitHub}
\noindent \textbf{Enlace:} \url{https://github.com/yguzmanm-sys/POO-2026-2S/blob/main/actividad%202/Actividad2_Ejercicio2_4_Figuras.py}

\subsection{Código fuente en Python}
\begin{lstlisting}[caption={Implementación de Figuras Geométricas (\texttt{Actividad2\_Ejercicio2\_4\_Figuras.py})}]
import math

class Circulo:
    def __init__(self, radio: float):
        self.radio = radio

    def calcular_area(self) -> float:
        return math.pi * (self.radio ** 2)

    def calcular_perimetro(self) -> float:
        return 2 * math.pi * self.radio

class Rectangulo:
    def __init__(self, base: float, altura: float):
        self.base = base
        self.altura = altura

    def calcular_area(self) -> float:
        return self.base * self.altura

    def calcular_perimetro(self) -> float:
        return 2 * (self.base + self.altura)

class Cuadrado:
    def __init__(self, longitud: float):
        self.longitud = longitud

    def calcular_area(self) -> float:
        return self.longitud ** 2

    def calcular_perimetro(self) -> float:
        return 4 * self.longitud

class TrianguloRectangulo:
    def __init__(self, base: float, altura: float):
        self.base = base
        self.altura = altura

    def calcular_area(self) -> float:
        return (self.base * self.altura) / 2

    def calcular_hipotenusa(self) -> float:
        return math.sqrt(self.base ** 2 + self.altura ** 2)

    def calcular_perimetro(self) -> float:
        return self.base + self.altura + self.calcular_hipotenusa()

    def determinar_tipo_triangulo(self) -> str:
        hipotenusa = self.calcular_hipotenusa()
        if math.isclose(self.base, self.altura) and math.isclose(self.base, hipotenusa):
            return "Es un triángulo Equilátero"
        elif not math.isclose(self.base, self.altura) and not math.isclose(self.base, hipotenusa) and not math.isclose(self.altura, hipotenusa):
            return "Es un triángulo Escaleno"
        else:
            return "Es un triángulo Isósceles"

    # Alias por compatibilidad
    determinar_tipo_triagunlo = determinar_tipo_triangulo

class Figuras_geometricas:
    @staticmethod
    def main():
        figura1 = Circulo(2.0)
        figura2 = Rectangulo(1, 2)
        figura3 = Cuadrado(3)
        figura4 = TrianguloRectangulo(3, 5)

        print(f"El área del círculo es = {figura1.calcular_area()}")
        print(f"El perímetro del círculo es = {figura1.calcular_perimetro()}")
        print(f"El área del rectángulo es = {figura2.calcular_area()}")
        print(f"El perímetro del rectángulo es = {figura2.calcular_perimetro()}")
        print(f"El área del cuadrado es = {figura3.calcular_area()}")
        print(f"El perímetro del cuadrado es = {figura3.calcular_perimetro()}")
        print(f"El área del triángulo rectángulo es = {figura4.calcular_area()}")
        print(f"El perímetro del triángulo rectángulo es = {figura4.calcular_perimetro()}")
        print(figura4.determinar_tipo_triangulo())

if __name__ == "__main__":
    Figuras_geometricas.main()
\end{lstlisting}

\subsection{Diagrama de Clases UML}
El diagrama ilustra las cuatro clases que representan las figuras geométricas independientes y la clase de prueba \texttt{Figuras\_geometricas} con una relación de uso (\textit{dependency/use}) hacia cada figura:

\begin{figure}[H]
    \centering
    \resizebox{\textwidth}{!}{
    \begin{tikzpicture}
        \node[uml, text width=3.6cm, anchor=north west] (circ) at (0,0) {\textbf{Circulo}
        \nodepart{second}\umlm radio: float
        \nodepart{third}\umlp \_\_init\_\_(radio: float)\\
        \umlp calcular\_area(): float\\
        \umlp calcular\_perimetro(): float};

        \node[uml, text width=3.6cm, anchor=north west] (rect) at ([xshift=0.4cm]circ.north east) {\textbf{Rectangulo}
        \nodepart{second}\umlm base: float\\
        \umlm altura: float
        \nodepart{third}\umlp \_\_init\_\_(base: float, altura: float)\\
        \umlp calcular\_area(): float\\
        \umlp calcular\_perimetro(): float};

        \node[uml, text width=3.6cm, anchor=north west] (cuad) at ([xshift=0.4cm]rect.north east) {\textbf{Cuadrado}
        \nodepart{second}\umlm longitud: float
        \nodepart{third}\umlp \_\_init\_\_(longitud: float)\\
        \umlp calcular\_area(): float\\
        \umlp calcular\_perimetro(): float};

        \node[uml, text width=4.5cm, anchor=north west] (tri) at ([xshift=0.4cm]cuad.north east) {\textbf{TrianguloRectangulo}
        \nodepart{second}\umlm base: float\\
        \umlm altura: float
        \nodepart{third}\umlp \_\_init\_\_(base: float, altura: float)\\
        \umlp calcular\_area(): float\\
        \umlp calcular\_hipotenusa(): float\\
        \umlp calcular\_perimetro(): float\\
        \umlp determinar\_tipo\_triangulo(): str};

        \node[umltwo, text width=5cm, anchor=north] (fig) at ($(rect.south)!0.5!(cuad.south)+(0,-2.3cm)$) {\textbf{Figuras\_geometricas}
        \nodepart{second}\umlp main(): None \quad \{static\}};

        \draw[dashed, -{Latex[length=2.5mm]}] (fig.north) -- node[near end, left, font=\scriptsize] {usa} (circ.south);
        \draw[dashed, -{Latex[length=2.5mm]}] (fig.north) -- (rect.south);
        \draw[dashed, -{Latex[length=2.5mm]}] (fig.north) -- (cuad.south);
        \draw[dashed, -{Latex[length=2.5mm]}] (fig.north) -- (tri.south);
    \end{tikzpicture}
    }
    \caption{Diagrama de clases UML vectorial para el Ejercicio 2.4 (Figuras Geométricas).}
    \label{fig:uml_ej2_4}
\end{figure}

\subsection{Resultado de ejecución}
\begin{verbatim}
El área del círculo es = 12.566370614359172
El perímetro del círculo es = 12.566370614359172
El área del rectángulo es = 2
El perímetro del rectángulo es = 6
El área del cuadrado es = 9
El perímetro del cuadrado es = 12
El área del triángulo rectángulo es = 7.5
El perímetro del triángulo rectángulo es = 13.8309518948453
Es un triángulo Escaleno
\end{verbatim}

\newpage

% =========================================================================
% EJERCICIO 2.5
% =========================================================================
\section{Ejercicio 2.5 — Métodos con Parámetros y Paso de Mensajes: Clase CuentaBancaria}
\textbf{Referencia en libro:} Capítulo 2, página 95 (y resolución en páginas 96--103).

\subsection{Enunciado del problema}
Se requiere un programa que modele una cuenta bancaria con los siguientes atributos:
\begin{itemize}
    \item Nombres del titular (\texttt{str}).
    \item Apellidos del titular (\texttt{str}).
    \item Número de cuenta bancaria (\texttt{int}).
    \item Tipo de cuenta (ahorros o corriente).
    \item Saldo de la cuenta (\texttt{float}, con valor inicial cero).
\end{itemize}

Operaciones requeridas en la cuenta bancaria:
\begin{itemize}
    \item Constructor que inicialice los atributos del titular y el número de cuenta.
    \item Método para imprimir los valores de los atributos.
    \item Método para consultar el saldo.
    \item Método para \textbf{consignar} un valor mayor que cero actualizando el saldo.
    \item Método para \textbf{retirar} un valor, validando que sea mayor a cero y no sobrepase el saldo actual.
    \item Métodos complementarios de interacción entre cuentas: \textbf{comparación} de cuentas y \textbf{transferencia} de fondos entre dos cuentas reutilizando las operaciones de retiro y consignación.
\end{itemize}

\subsection{URL del repositorio en GitHub}
\noindent \textbf{Enlace:} \url{https://github.com/yguzmanm-sys/POO-2026-2S/blob/main/actividad%202/Actividad2_Ejercicio2_5_CuentaBancaria.py}

\subsection{Código fuente en Python}
\begin{lstlisting}[caption={Implementación de la clase CuentaBancaria (\texttt{Actividad2\_Ejercicio2\_5\_CuentaBancaria.py})}]
class CuentaBancaria:
    def __init__(self, nombresTitular: str, apellidosTitular: str,
                 numeroCuenta: int, tipoCuenta: str):
        self.nombresTitular = nombresTitular
        self.apellidosTitular = apellidosTitular
        self.numeroCuenta = numeroCuenta
        self.tipoCuenta = tipoCuenta
        self.saldo = 0.0

    def imprimir(self):
        print("Nombres del titular =", self.nombresTitular)
        print("Apellidos del titular =", self.apellidosTitular)
        print("Número de cuenta =", self.numeroCuenta)
        print("Tipo de cuenta =", self.tipoCuenta)
        print("Saldo =", self.saldo)

    def consultarSaldo(self):
        print("El saldo actual es =", self.saldo)
        return self.saldo

    def consignar(self, valor: float) -> bool:
        if valor > 0:
            self.saldo += valor
            print(f"Se ha consignado ${valor}. Nuevo saldo: ${self.saldo}")
            return True
        print("El valor a consignar debe ser mayor que cero.")
        return False

    def retirar(self, valor: float) -> bool:
        if 0 < valor <= self.saldo:
            self.saldo -= valor
            print(f"Se ha retirado ${valor}. Nuevo saldo: ${self.saldo}")
            return True
        print("El valor a retirar debe ser menor o igual al saldo actual y mayor que cero.")
        return False

    def compararCuenta(self, cuenta: 'CuentaBancaria') -> bool:
        return (self.numeroCuenta == cuenta.numeroCuenta
                and self.nombresTitular == cuenta.nombresTitular
                and self.apellidosTitular == cuenta.apellidosTitular
                and self.tipoCuenta == cuenta.tipoCuenta)

    def transferencia(self, cuenta: 'CuentaBancaria', valor: float) -> bool:
        if self.retirar(valor):
            cuenta.consignar(valor)
            print(f"Se ha transferido ${valor} de la cuenta {self.numeroCuenta} a {cuenta.numeroCuenta}")
            return True
        print("No se pudo realizar la transferencia.")
        return False

if __name__ == "__main__":
    cuenta1 = CuentaBancaria("Pedro", "Pérez", 123456, "AHORROS")
    cuenta2 = CuentaBancaria("Juan", "López", 654321, "CORRIENTE")

    cuenta1.imprimir()
    cuenta1.consignar(500000)
    cuenta1.transferencia(cuenta2, 200000)
    cuenta1.consultarSaldo()
    cuenta2.consultarSaldo()
\end{lstlisting}

\subsection{Diagrama de Clases UML}
El diagrama modela la clase \texttt{CuentaBancaria} y la enumeración \texttt{tipoCuenta}:

\begin{figure}[H]
    \centering
    \begin{tikzpicture}
        \node[uml, text width=8.5cm] (cuenta) {\textbf{CuentaBancaria}
        \nodepart{second}\umlm nombresTitular: str\\
        \umlm apellidosTitular: str\\
        \umlm numeroCuenta: int\\
        \umlm tipoCuenta: str\\
        \umlm saldo: float = 0.0
        \nodepart{third}\umlp \_\_init\_\_(nombresTitular: str, apellidosTitular: str, numeroCuenta: int, tipoCuenta: str)\\
        \umlp imprimir(): None\\
        \umlp consultarSaldo(): float\\
        \umlp consignar(valor: float): bool\\
        \umlp retirar(valor: float): bool\\
        \umlp compararCuenta(cuenta: CuentaBancaria): bool\\
        \umlp transferencia(cuenta: CuentaBancaria, valor: float): bool};

        \node[umltwo, text width=3.4cm, right=1.6cm of cuenta] (tcuenta) {\estereo{enumeration}\\\textbf{tipoCuenta}
        \nodepart{second}AHORROS\\
        CORRIENTE};

        \draw[dashed, -{Latex[length=2.5mm]}] (cuenta.east) -- node[above, font=\scriptsize] {valores} (tcuenta.west);
    \end{tikzpicture}
    \caption{Diagrama de clases UML vectorial para el Ejercicio 2.5 (Clase CuentaBancaria).}
    \label{fig:uml_ej2_5}
\end{figure}

\subsection{Resultado de ejecución}
\begin{verbatim}
Nombres del titular = Pedro
Apellidos del titular = Pérez
Número de cuenta = 123456
Tipo de cuenta = AHORROS
Saldo = 0.0
Se ha consignado $500000. Nuevo saldo: $500000.0
Se ha retirado $200000. Nuevo saldo: $300000.0
Se ha consignado $200000. Nuevo saldo: $200000.0
Se ha transferido $200000 de la cuenta 123456 a 654321
El saldo actual es = 300000.0
El saldo actual es = 200000.0
\end{verbatim}

\newpage

% =========================================================================
% CONCLUSIONES
% =========================================================================
\section{Conclusiones}
\begin{itemize}[leftmargin=*]
    \item La programación orientada a objetos en Python permite un modelado conceptual intuitivo y flexible mediante clases, constructores y métodos, manteniendo una correspondencia directa con los diagramas estáticos de UML.
    \item El uso de tipos enumerados (\texttt{Enum}) aporta robustez y claridad al código, limitando el rango de valores permitidos a constantes predefinidas semánticamente coherentes con el dominio del problema.
    \item Los mecanismos de encapsulamiento mediante métodos de acceso (\textit{getters} y \textit{setters}) permiten controlar la mutabilidad y validar las transiciones de estado de los objetos, evitando inconsistencias como velocidades negativas o desbordamiento de límites operativos.
    \item El paso de objetos como parámetros entre métodos permite implementar patrones de interacción y delegación de responsabilidades de alta cohesión, como se evidenció en la transferencia de saldos entre cuentas bancarias.
\end{itemize}

\end{document}
